\scanpage{77}
I) \unsure{Condition|Conditions} de finitude, sur corps de base $k$

\uline{Prop.} Si $G_0$ est fini, alors \deleted{\unsure{les}} $\beta_0,\beta_1$ sont algébriques sur le corps premier. L'inverse est vrai si \deleted{$p=$} $\operatorname{car}.k>0$.

% slanted
\margin{valable en dim quelconque}

\deleted{\uline{\unsure{Raisons|Racines}} \unsure{de ce que}}

Introduisons \deleted{\unsure{les val propres}}\inserted{des solutions} (dans une clôture alg. $\bar k$ de $k$) $\zeta_i$ $(i=0,1)$ de l'équation aux val. propres

\[
 \zeta_i^2-\alpha_i\zeta_i+1=0\qquad\text{i.e.}\quad \zeta_i+\zeta_i^{-1}=\alpha_i
\]

Donc pour $i=0,1$ fixé, $\zeta_i$ est déterminé mod passage à $\zeta_i^{-1}$.

\[
 \text{Soit}\qquad\unsure{\nu'_i|\nu_i}=\text{ordre de }\zeta_i\qquad i=0,1
\]

Donc $\nu_i$ est tjrs premier à $p$, sauf si $\nu'_i=1$ i.e. $\zeta_i=1$, i.e. $\alpha_i=2$.

\unsure{Rappelons}

\uline{Prop} a) Si $\nu'_i\ne\deleted{\unsure{1}}\inserted{1,2}$ \deleted{\unsure{alors}} \inserted{i.e. $\alpha_i\ne\pm2$ i.e. $\beta_i(4-\beta_i)\ne0$,} alors $\nu_i=$ ordre de $\rho_i$

(\uline{NB} $\rho_0=\rho_f,\rho_1=\rho_s$)

b) Si \deleted{\unsure{$\nu'_i=2$}}\inserted{i.e. $\alpha_i=-2$, i.e. $\deleted{\unsure{\alpha}}\beta_i=0$ ($\nu'_i=2$, ou $\nu'_i=1$ en car $2$)} alors $\rho_i$ d'ordre $2p$

c) \deleted{Si \unsure{$\nu'_i=1$} en car \unsure{$p$} $\alpha_i=2\ne-2$, \unsure{$\rho_i$} alors $\left\{\begin{smallmatrix}\nu_i=p&\text{si }p\ne2\\\nu_i=\unsure{1}&\text{si }p=2\end{smallmatrix}\right.$}
Si $\alpha_i=2\ne-2$ (i.e. $\nu'_i=1$ en car. $\ne2$) alors $\nu_i=p$

[\,\uline{NB} Dans b) et c), on convient que $p=+\infty$ dans cas de car. nulle\,]

\uline{Théorème} Supposons $k$ de car. $0$. Pour que $G_0$ soit fini, il f et s. que l'on soit dans l'un des \deleted{quatre}\inserted{six} cas suivants

\deleted{1) $\nu_0=3,\nu_1=3$ (tétraèdre)}
\deleted{2) $\nu_0=3,\nu_1=4$ (octaèdre) \quad 2') $\nu_0=4,\nu_1=3$ (cube)}
\deleted{3) $\nu_0=3,\nu_1=5$ (icosaèdre) \quad 3') $\nu_0=5,\nu_1=3$ (dodécaèdre)}
\deleted{\unsure{2) $\nu_0=3,\nu_1=$}}

% slanted
\margin{on appellera ces cas les cas pythagoriciens (\unsure{sur base quelconque|cas base quelconque})}

\[
 \mathrm{I})\quad\nu_0=3,\ \nu_1=3\quad(\text{tétraèdre})\qquad
 \text{i.e. }\alpha_0=\alpha_1=-1,\quad\text{i.e. }\beta_0=\beta_1=+1\ \searrow
\]

\[
 \mathrm{II})\quad\nu_0=3,\ \nu_1=4\quad(\text{octaèdre})\qquad
 \text{i.e. }\alpha_0=-1,\alpha_1=0\quad\text{i.e. }\beta_0=1,\beta_1=2
\]

\[
 \mathrm{III})\quad\nu_0=3,\ \nu_1=5\quad(\text{icosaèdre})\quad
 \text{i.e. }\alpha_0=-1,\ \alpha_1\deleted{\unsure{=\alpha}}\text{ solution de l'équation }
 \boxed{\alpha_1^2+\alpha_1-1=0}
\]

\[
 \mathrm{II}')\quad\nu_0=4,\ \nu_1=3\quad(\text{cube})\qquad
 \text{i.e. }\alpha_0=0,\alpha_1=-1\quad\text{i.e. }\beta_0=2,\beta_1=1
\]

\[
 \mathrm{III}')\quad\nu_0=5,\ \nu_1=3\quad(\text{dodécaèdre})\quad
 \text{i.e. }\alpha_0\text{ solution de }\boxed{\alpha_0^2+\alpha_0-1=0},\quad
 \alpha_1=-1\ \text{i.e. }\beta_1=1
\]

\[
 \mathrm{IV})\quad\nu_0=\nu_1=5,\ \uline{\text{et}}\ \alpha_0\ne\alpha_1\
 \overset{\text{(pentagramme)}}{\vphantom{I}}\quad\text{i.e. }\alpha_0,\alpha_1
 \text{ solutions }\uline{\text{distinctes}}\text{ de }\boxed{\alpha^2+\alpha-1=0}
\]
% archive sheet number: 76
